\section{Билет 9. Представление гармонического колебания с помощью векторной диаграммы. Сложение гармонических колебаний одного направления. Амплитуда результирующего колебаний.}

\begin{definition}
Гармонические колебания удобно представлять графически с помощью \textbf{векторных диаграмм}. Построим вектор $\vec{A}$ длиной, равной амплитуде колебаний $A$, под углом $\varphi$ к горизонтальной оси $Ox$. Если привести этот вектор во вращение с угловой скоростью $\omega_0$, то проекция его конца на ось $Ox$ будет изменяться по закону:
\begin{equation}
    x(t) = A \cos(\omega_0 t + \varphi). \label{eq:vector_diag}
\end{equation}
Таким образом, проекция конца вращающегося вектора совершает гармонические колебания с амплитудой $A$, круговой частотой $\omega_0$ и начальной фазой $\varphi$.
\end{definition}

\begin{theorem}[Сложение колебаний одного направления]
Рассмотрим сложение двух гармонических колебаний одинакового направления и одинаковой частоты $\omega_0$:
\begin{equation}
    x_1(t) = A_1 \cos(\omega_0 t + \varphi_1), \quad x_2(t) = A_2 \cos(\omega_0 t + \varphi_2).
\end{equation}
Результирующее смещение равно сумме смещений: $x(t) = x_1(t) + x_2(t)$. В методе векторных диаграмм это соответствует сложению векторов $\vec{A}_1$ и $\vec{A}_2$:
\begin{equation}
    \vec{A} = \vec{A}_1 + \vec{A}_2.
\end{equation}
Поскольку оба вектора вращаются с одной и той же угловой скоростью $\omega_0$, результирующий вектор $\vec{A}$ также вращается с частотой $\omega_0$. Следовательно, суммарное движение остаётся гармоническим колебанием той же частоты:
\begin{equation}
    x(t) = A \cos(\omega_0 t + \varphi). \label{eq:res_osc}
\end{equation}
Амплитуда $A$ результирующего колебания определяется по теореме косинусов:
\begin{equation}
    A = \sqrt{A_1^2 + A_2^2 + 2A_1 A_2 \cos(\varphi_2 - \varphi_1)}. \label{eq:amplitude_sum}
\end{equation}
\end{theorem}

\begin{property}
Разность фаз складываемых колебаний $\Delta\varphi = \varphi_2 - \varphi_1$ не зависит от времени ($\Delta\varphi = \text{const}$). Такие колебания называются \textbf{когерентными}. Амплитуда результирующего колебания зависит от разности фаз:
\begin{enumerate}
    \item Если $\Delta\varphi = 2m\pi$ ($m=0, 1, 2, \dots$), колебания \textbf{синфазны}. Амплитуда максимальна: $A = A_1 + A_2$.
    \item Если $\Delta\varphi = (2m+1)\pi$ ($m=0, 1, 2, \dots$), колебания находятся в \textbf{противофазе}. Амплитуда минимальна: $A = |A_1 - A_2|$.
\end{enumerate}
\end{property}

\begin{remark}
Начальная фаза результирующего колебания $\varphi$ определяется из геометрического сложения проекций:
\begin{equation}
    \tan\varphi = \frac{A_1 \sin\varphi_1 + A_2 \sin\varphi_2}{A_1 \cos\varphi_1 + A_2 \cos\varphi_2}.
\end{equation}
Таким образом, сложение когерентных колебаний одного направления и одинаковой частоты всегда даёт новое гармоническое колебание с той же частотой, но с амплитудой и фазой, зависящими от параметров слагаемых колебаний.
\end{remark}
